% !TeX root = ../../Book3.tex
% 纲要标题：### E.3 特征集与有效消去的代表方法
% 纲要位置：计划纲要.md，第 2369 行。

\section{特征集与有效消去的代表方法}
\label{b3:sec:E03}

% TODO: 按以下纲要要点撰写本节。

% 纲要内容（仅作写作计划，尚非教材正文）：
%
% * 导数排序、导变量、初式、分离子及伪约化；特征集不等于普通理想的一组生成元
% * Rosenfeld–Gröbner 的正则微分链与饱和分解概览；不能自动称为素分解或唯一不可约分解
% * 只给一个可复核的消去例子，保留分母为零、初式为零和分离子为零的分支；完整终止证明、界和实现另立专题
%

\subsection{导数排序、导变量、初式、分离子与伪约化}
\label{b3:subsec:E03:01}

待撰写。

\subsection{Rosenfeld–Gröbner 正则微分链与饱和分解}
\label{b3:subsec:E03:02}

待撰写。

\subsection{消去算例与分母、初式、分离子的退化分支}
\label{b3:subsec:E03:03}

待撰写。
